\section{Das Axiomsystem der Wortstrukturen}

Die Axiomatik beschreibt, welche Eigenschaften eine Menge von Wörtern
mit einem Anfangswort und einer Anfügungsfunktion besitzt. Dabei haben
die Symbole folgende Bedeutung:

\begingroup
\renewcommand{\arraystretch}{1.15}
\noindent\begin{tabular}{@{}p{.14\linewidth}p{.81\linewidth}@{}}
\textbf{Symbol} & \textbf{Bedeutung} \\
\(A\) & Das Alphabet, also die Menge der zugelassenen Buchstaben. \\
\(W\) & Die Trägermenge der Struktur; ihre Elemente werden als ganze Wörter aufgefasst. \\
\(e\) & Das ausgezeichnete Anfangswort in \(W\); es übernimmt die Rolle des leeren Wortes. \\
\(s\) & Die Anfügungsfunktion: \(s(u,a)\) entsteht aus dem Wort \(u\), indem rechts der Buchstabe \(a\) angefügt wird. \\
\(u,v\) & Beliebige Wörter, also Elemente von \(W\). \\
\(a,b\) & Beliebige Buchstaben, also Elemente von \(A\). \\
\(U\) & Eine beliebige Teilmenge von \(W\), an der die Minimalitätsbedingung geprüft wird.
\end{tabular}\par
\endgroup

Da \(s\) auf \(W\times A\) definiert ist, bedeutet \(s(u,a)\) den
Funktionswert am geordneten Paar \((u,a)\). Die Funktion hat somit zwei
Eingaben: ein bisheriges Wort und einen einzelnen Buchstaben.
Beispielsweise stellt \(s(e,a)\) das aus \(a\) bestehende Wort dar;
\(s(s(e,a),b)\) stellt das Wort mit den aufeinanderfolgenden Buchstaben
\(a,b\) dar. Die Buchstaben dürfen dabei gleich sein.

Im zuvor aufgebauten Anfügungsmodell werden die Symbole durch
\[
 W=\FiniteWords A,\qquad e=\EmptyWord A,\qquad s=\sigma_A
\]
belegt. Nach \FormulaRefAuto{EarlyWordSuccessorEquation}[thm] ist dort
für \(u\in\FiniteWords A\) und \(a\in A\)
\[
 \sigma_A(u,a)=u\cup\{(\WordLength u,a)\}.
\]
Der neue Buchstabe erhält also die nächste freie Position. In einer
beliebigen Wortstruktur sind \(W,e,s\) dagegen zunächst die gewählten
Strukturdaten; ihre konkrete Darstellung ist durch die Axiomatik nicht
festgelegt. Insbesondere muss \(e\) dort nicht die leere Menge sein.

Für eine Menge \(A\) heißt \((W,e,s)\) eine \emph{freie Wortstruktur über
\(A\)}, wenn die folgenden Bedingungen gelten.

\Needspace{19\baselineskip}
\FormulaDefDeltaK[Freie Wortstruktur]{%
  \mathsf{WortStr}_A(W,e,s)
}{WordStructureDef}{%
  \DeltaRow{Mengen}{A\dsep W\dsep e\dsep U\dsep u\dsep v\dsep a\dsep b}
  \DeltaRow{Funktionen}{s}
  \DeltaRow{\textbf{Strukturaxiome}}{}
  \DeltaRow{W0: Typisierung}{e\in W\land s\colon W\times A\to W}
  \DeltaRow{W1: Anfang}{\forall u\in W\,\forall a\in A\;s(u,a)\ne e}
  \DeltaRow{W2: Eindeutigkeit}{%
    \begin{aligned}[t]
      &\forall u,v\in W\,\forall a,b\in A\\[-2pt]
      &\quad(s(u,a)=s(v,b)\rightarrow u=v\land a=b)
    \end{aligned}}
  \DeltaRow{W3: Minimalität}{%
    \begin{aligned}[t]
      &\forall U\subseteq W\,\Bigl(\bigl(e\in U\land{}\\[-2pt]
      &\quad(\forall u\in U\,\forall a\in A\;s(u,a)\in U)\bigr)
        \rightarrow U=W\Bigr)
    \end{aligned}}
  \DeltaRow{\textbf{Neues Symbol}}{}
  \DeltaPrem{Freie Wortstruktur}{\mathsf{WortStr}_A(W,e,s)}
}

Die bei den einzelnen Axiomen genannten Herkunftssätze beweisen die
jeweilige Eigenschaft für das Anfügungsmodell. Für beliebige Daten
\((W,e,s)\) werden diese Eigenschaften in
\FormulaRefAuto{WordStructureDef}[def] gefordert. Unter der Annahme
\(\mathsf{WortStr}_A(W,e,s)\) geben die folgenden nummerierten Axiome
jeweils die entsprechende Bedingung dieser Definition wieder.

\noindent\begin{minipage}{\linewidth}
\FormulaAxiomDeltaK[Typisierung einer Wortstruktur]{%
  \mathsf{WortStr}_A(W,e,s)\vdash e\in W\land s\colon W\times A\to W
}{WordStructureTypingAxiom}{%
  \DeltaRow{Mengen}{A\dsep W\dsep e}\DeltaRow{Funktionen}{s}}

\smallskip\noindent\textit{Zu W0.} Das Anfangswort gehört zu \(W\), und
die Anfügung ist für jedes Paar aus einem Wort und einem Buchstaben
definiert und liefert wieder ein Wort. Für
\((\FiniteWords A,\EmptyWord A,\sigma_A)\) ist dies genau
\FormulaRefAuto{EarlyWordSuccessorTyping}[thm].
\end{minipage}\par

\noindent\begin{minipage}{\linewidth}
\FormulaAxiomDeltaK[Das Anfangswort ist kein Nachfolger]{%
  \mathsf{WortStr}_A(W,e,s)\dsep u\in W\dsep a\in A\vdash s(u,a)\ne e
}{WordStructureZeroAxiom}{%
  \DeltaRow{Mengen}{A\dsep W\dsep e\dsep u\dsep a}\DeltaRow{Funktionen}{s}}

\smallskip\noindent\textit{Zu W1.} Durch das Anfügen eines Buchstabens
entsteht niemals das Anfangswort. Im Anfügungsmodell ist diese Aussage
durch \FormulaRefAuto{EarlyWordSuccessorNonempty}[thm] bewiesen.
\end{minipage}\par

\noindent\begin{minipage}{\linewidth}
\FormulaAxiomDeltaK[Gemeinsame Injektivität der Wortanfügung]{%
  \begin{aligned}[t]
    &\mathsf{WortStr}_A(W,e,s)\dsep u,v\in W\dsep a,b\in A\dsep{}\\[-2pt]
    &s(u,a)=s(v,b)\vdash u=v\land a=b
  \end{aligned}
}{WordStructureInjectivityAxiom}{%
  \DeltaRow{Mengen}{A\dsep W\dsep e\dsep u\dsep v\dsep a\dsep b}
  \DeltaRow{Funktionen}{s}}

\smallskip\noindent\textit{Zu W2.} Gleiche angefügte Wörter haben sowohl
denselben bisherigen Wortteil als auch denselben letzten Buchstaben.
Die beiden Eingaben lassen sich also eindeutig zurückgewinnen.
Injektivität nur bei festgehaltenem Buchstaben würde dafür nicht
genügen. Der konkrete Nachweis ist
\FormulaRefAuto{EarlyWordSuccessorInjective}[thm].
\end{minipage}\par

\noindent\begin{minipage}{\linewidth}
\FormulaAxiomDeltaK[Minimalität einer Wortstruktur]{%
  \begin{aligned}[t]
    &\mathsf{WortStr}_A(W,e,s)\dsep U\subseteq W\dsep e\in U\dsep{}\\[-2pt]
    &\forall u\in U\,\forall a\in A\;s(u,a)\in U\vdash U=W
  \end{aligned}
}{WordStructureMinimalityAxiom}{%
  \DeltaRow{Mengen}{A\dsep W\dsep e\dsep U\dsep u\dsep a}
  \DeltaRow{Funktionen}{s}}

\smallskip\noindent\textit{Zu W3.} Enthält eine Teilmenge \(U\) das
Anfangswort und mit jedem ihrer Wörter auch alle Anfügungen, so umfasst
sie bereits ganz \(W\). Damit gibt es keine echte Teilmenge von \(W\),
die unter diesen Erzeugungsregeln abgeschlossen ist und den Anfang
enthält. Im Anfügungsmodell liefert
\FormulaRefAuto{EarlyWordModelMinimality}[thm] genau diese Aussage.
\end{minipage}\par
